Caching the work of a \emph{declared} build action is routine: the action states its inputs, so a
build system or a compiler cache can key on them. Caching an arbitrary, already-written command is
harder, because the command does not state its inputs, and replaying its output without knowing
them is a stale-result hazard. This paper studies the other side of that trade-off: what a
deliberately conservative automatic command cache costs, and what it buys.

We describe \tracejit{}, a whole-command cache for unmodified Linux~x86\_64 processes. It traces
one execution with \texttt{ptrace}, classifies every kernel observation, compiles the observations
into synchronous guards, and reuses the recorded standard streams, exit status and output files
only when every guard still passes. Anything it cannot re-check --- the clock, entropy, process
identity, the network, an unmodelled syscall --- is refused rather than assumed. Two execution
tiers are kept strictly apart: \emph{native} execution, and an explicit, opt-in \emph{controlled}
contract that supplies a small set of runtime values reproducibly.

On a frozen 59-workload census, native reuse eligibility is $2/49$ ($4.1\%$); under the audited
\texttt{allocator-entropy-v1} contract it is $12/49$ ($24.5\%$), and under
\texttt{compiler-runtime-v1} $19/49$ ($38.8\%$). Across all cells, 28 of 162 attempted rows reused
and none produced an incorrect observation. On a same-host comparison against \firebuild{}
0.8.2 and 0.8.6.2, \tracejit{} is dramatically better on one workload class and worse on others:
a 25-invocation \texttt{bzip2} session costs 545\,ms against 2847\,ms natively ($5.22\times$,
$80.8\%$ less cumulative time) while both \firebuild{} versions decline to reuse it at all;
\firebuild{} is faster on file-output \texttt{gzip}, on a tuned \texttt{sed}, and on hash-only work
below 10\,MiB, and \texttt{ccache} is far faster than either on compilation.

These are negative results as much as positive ones, and we report them as such. Reuse requires an
audited libc build, so most ordinary programs are ineligible on an unaudited one. Verified replay
of a 100\,MiB artifact costs tens to hundreds of milliseconds, dominated by hashing and output
materialisation, so cheap commands lose. Enforced execution --- seccomp plus Landlock, a state we
call \proven{} --- completes a full create/reuse/invalidate lifecycle for read-only commands on a
capable host, but cannot currently promote the artifact producers that are the product's stated
niche. Most importantly, a writer that mutates an input \emph{after} guard validation can still
cause stale output to be served; we reproduce that counterexample and state it as an operating
assumption rather than a guarantee.

The contribution is not a new caching primitive. It is a measured account of conservative
automatic reuse: a contract discipline in which refusal is the default, the price of that
discipline in eligibility, latency and applicability, and the exact set of correctness boundaries
that remain open.
